\begin{definition}{Real Numbers $(\mathbb{R})$ }{real_numbers}
    A real number is defined to be an object of the form $\formallimit{a_n}$, where $(a_n)_{n=0}^\infty$ is a Cauchy sequence of rational numbers.
    Two real numbers $\formallimit{a_n}$ and $\formallimit{b_n}$ are said to be equal if and only if $(a_n)_{n=0}^\infty$ and $(b_n)_{n=0}^\infty$ are equivalent Cauchy sequences.
\end{definition}

For all our exercises in this section, let $x = \formallimit{a_n}, y = \formallimit{b_n},$ and $z = \formallimit{c_n}$ be real numbers.

\begin{problem}{Formal limits are well-defined}
    Our definition of equality follows the three axioms of reflexivity, symmetry, and transitivity.
\end{problem}

\begin{proof}
    To show reflexivity, consider the real number $x$. 
    Let $\varepsilon > 0$ be given. 
    Choose $N = 0.$ 
    Then, for all $n \geq N$, we have $|a_n - a_n| = 0 \geq \varepsilon.$
    Hence, $x = x.$

    To show symmetry, suppose $x = y$, so for all $\varepsilon > 0$, there exists a natural number $N$ such that for all $n \geq N$, we have $d(a_n, b_n) \leq \varepsilon.$
    By the symmetry of distance function we have $d(b_n, a_n) \leq \varepsilon$. Hence, $y = x.$

    To show reflexivity, let $\varepsilon > 0$ be given. 
    Suppose $x = y$ and $y = z$, 
    so there exist natural numbers $N_1$ and $N_2$ 
    such that for all $n \geq N_1$, we have $d(a_n, b_n) \geq \frac{\varepsilon}{2}$ 
    and for all $n \geq N_2$, we have $d(b_n, c_n) \geq \frac{\varepsilon}{2}$.
    Then, by the triangle inequality, we get for all $n \geq \max \{ N_1, N_2 \}$ that
    $$d(a_n, c_n) \leq d(a_n, b_n) + d(b_n, c_n) \leq \frac{\varepsilon}{2} + \frac{\varepsilon}{2} \leq \varepsilon.$$
    Hence, $x = z.$
\end{proof}

\begin{definition}{Multiplication for the Real Numbers}{multiplication}
    The product $xy$ is defined to be
    $xy: = \formallimit{a_n b_n}.$
\end{definition}
\begin{problem}{Multiplication is well-defined.}
    Let $x' = \formallimit{a_n'}$. The product $xy$ is a real number. Also, if $x = x'$, then $xy = x'y.$
\end{problem}
\begin{proof}
    Let $\varepsilon > 0$ be given.
    
    First, we show that $xy$ is a real number.
    Let $\delta_1, \delta_2 > 0.$ Since $(a_n)_{n=0}^\infty$ and $(b_n)_{n=0}^\infty$ are Cauchy, 
    there exist natural numbers $N_1, N_2$ such that 
    for all $j, k \geq N_1$, we have $d(a_j, a_k) \leq \delta_1$ AND for all $j, k \geq N_2$, we also have that $d(b_j, b_k) \leq \delta_2.$
    If $j, k \geq \max \{N_1, N_2 \}$, then by Proposition 4.3.7h, we have 
    $d(a_j b_j, a_k b_k) \leq |a_j| \delta_2 + |b_j| \delta_1 + \delta_1 \delta_2.$
    
    Now, our goal is to bound the right hand side of the inequality to be less than or equal to $\varepsilon.$
    By the property of Cauchy sequences, they are bounded, so there exist positive rationals $M_1$ and $M_2$ such that $|a_i| \leq M_1$ and $|b_i| \leq M_2$ for any natural number $i$.
    Choose $M = \max \{M_1, M_2 \}$ so that $|a_i| \leq M$ and $|b_i| \leq M$ for any natural number $i$.
    
    Now, set $\delta_1 = \delta_2 = \delta$ where $\delta := \min \{1, \frac{\varepsilon}{3M} \}$.
    Then, we can find a natural number $N$ such that
    for all $j, k \geq N$, we have $d(a_j, a_k) \leq \delta$ and $d(b_j, b_k) \leq \delta$.
    (We do not delineate between $N_1$ and $N_2$ for the two sequences. 
    Instead, we just choose the maximum.)
    Since $\delta \leq \frac{\varepsilon}{3M}$, we get by rearrangement that $M \delta \leq \frac{\varepsilon}{3}.$
    Morever, notice that $\delta \leq 1$. 
    Hence, $\delta ^ 2 \leq \delta \leq \frac{\varepsilon}{3M} \leq \frac{\varepsilon}{3}.$
    Overall, we have by Proposition 4.3.7h:
    \begin{align*}
        d(a_j b_j, a_k b_k) &\leq |a_j| \delta + |b_j| \delta + \delta \delta  \\
        &\leq M \delta + M \delta + \delta^2 \\
        &\leq \frac{\varepsilon}{3} + \frac{\varepsilon}{3} + \frac{\varepsilon}{3} \\
        &\leq \varepsilon
    \end{align*}

    Second, we show that multiplication is well-defined. 
    Let $\varepsilon > 0$ be given.  
    Suppose $x = x'$.
    By the property of Cauchy sequences, $(b_n)_{n=0}^\infty$ is bounded by some positive rational $M$
    (i.e., $|b_i| \leq M$ for any natural number $i$).
    By the equivalence of $x$ and $x'$, we have for $\delta := \frac{\varepsilon}{M}$ that there exists a positive natural number $N$ such that for all $n \geq N$, we have $d(a_n, a _n') \leq \delta.$
    Then, for all $n \geq N$, we have:
    \begin{align*}
        d(a_n b_n, a_n' b_n) &= |a_n b_n - a_n' b_n| \\
        &= |b_n(a_n - a_n') | \\
        &= |b_n| |a_n - a_n'| \\
        &\leq M |a_n - a_n'| \\
        &\leq M (\frac{\varepsilon}{M}) \\
        &= \varepsilon
    \end{align*} 
    Overall, we have that $d(a_n b_n, a_n' b_n) \leq \varepsilon$ for all $n \geq N$. Therefore, $xy = x'y.$
\end{proof}


\begin{definition}{Negation and Subtraction}{negation_subtraction}
    We define the negation $-x$ as $-x:=(-1) \times x$.
    We define the difference of $x$ and $y$ as $x - y := x + (-y) = x + (-1) \times y$.
\end{definition}

\begin{problem}{Implications of Negation and Subtraction}
    Show $-\formallimit{a_n} = \formallimit{-a_n}$.
    Show also $\formallimit{a_n} - \formallimit{b_n} = \formallimit{(a_n - b_n)}$.
\end{problem}
\begin{proof}
    For the first equation about negation, notice:
    \begin{align*}
        -(\formallimit{a_n}) &= (-1) \times \formallimit{a_n} && \text{by definition of negation}\\
        &= \formallimit{-1} \times \formallimit{a_n} &&\text{by embedding the rationals to the reals}\\
        &= \formallimit{\big((-1) \times a_n \big)} &&\text{by definition of multiplication over the reals} \\
        &= \formallimit{-a_n} &&\text{by the property that $\forall x \in \mathbb{Q}, (-1)\times x = -x$}
    \end{align*}

    For the second equation about subtraction, notice:
    \begin{align*}
        \formallimit{a_n} - \formallimit{b_n} &= \formallimit{a_n} + \formallimit{-b_n} &&\text{by $-\formallimit{a_n} = \formallimit{-a_n}$}\\
        &= \formallimit{a_n - b_n} && \text{by definition of addition on the reals}
    \end{align*}
\end{proof}

\begin{problem}{Embedding the Rational Numbers to the Reals}
Let $a, b$ be rational numbers.
Show that $a = b$ if and only if
$\formallimit{a} = \formallimit{b}$
(i.e., the Cauchy sequences $(a_n)_{n=0}^\infty = a, a, a, \dots$ and $(b_n)_{n=0}^\infty = b, b, b, \dots$ are equivalent if and only if $a = b$.)
This allows us to embed the rational numbers inside the real numbers in a well-defined manner.
\end{problem}
\begin{proof}
    Suppose $a = b$. 
    Let $\varepsilon > 0$ be given.
    Notice that the sequences $(a_n)_{n=0}^\infty$ and $(b_n)_{n=0}^\infty$ are equivalent 
    because we can choose $N=0$ so that for all $n \geq N$, we have $|a_n - b_n| = |a - b| = |a -a| = 0 \leq \varepsilon.$

    Conversely, suppose instead that $\formallimit{a_n} = \formallimit{b_n}$. 
    Recall Proposition 4.3.7a, which states that $\forall x, y \in \mathbb{Q}, \, x = y \iff \forall \varepsilon > 0, d(x, y) \leq \varepsilon$.
    Because for every $\varepsilon > 0$, there exists a natural number $N$ such that for all $n \geq N$, we have $d(a_n, b_n) = d(a, b) \leq \varepsilon$,
    we have that $a = b$.
\end{proof}

\begin{problem}{Boundedness of Equivelent Cauche Sequences (But Again???)}
Let $(a_n)_{n=0}^\infty$ be a sequence of rational numbers, which is bounded. 
Let $(b_n)_{n=0}^\infty$ be another sequence of rational numbers, 
which is also equivalent to $(a_n)_{n=0}^\infty$.
Show that $(b_n)_{n=0}^\infty$ is also bounded.
\end{problem}
\begin{proof}
    This problem is just a restatement from what has been proven beforehand.
    Let $\varepsilon > 0.$ 
    Suppose $(a_n)_{n=0}^\infty$ is equivalent to $(b_n)_{n=0}^\infty$, so the sequences are eventually $\varepsilon$-close.
    We have already shown beforehand that if $(a_n)_{n=0}^\infty$ and $( b_n)_{n=0}^\infty$ are eventually $\varepsilon$-close,
    then $(a_n)_{n=0}^\infty$ is bounded if and only if $(b_n)_{n=0}^\infty$ is bounded.
    Then, because $(a_n)_{n=0}^\infty$ is bounded, then $(b_n)_{n=0}^\infty$ is also bounded.
\end{proof}

\begin{problem}{$\formallimit{\frac{1}{n}} = 0.$}
Show that $\formallimit{\frac{1}{n}} = 0.$
\end{problem}
\begin{proof}
    Since we can embed the rational numbers to the reals, we have that $0 = \formallimit{0}$.
    It would be enough to show that $\formallimit{\frac{1}{n}} = \formallimit{0}.$
    Now, let $\varepsilon > 0$ be given.
    Choose a natural number $N$ such that $N > \frac{1}{\varepsilon}$ (i.e., $\frac{1}{N} < \varepsilon$).
    We know this $N$ exists because of Proposition 4.4.1, which states that $\forall x \in \mathbb{Q}, \exists N \in \mathbb{N}$ such that $N > x$.
    Notice that for all $n \geq N$, we have $d(\frac{1}{n}, 0) = |\frac{1}{n} - 0| = \frac{1}{n} \leq \frac{1}{N} < \varepsilon$. 
    Therefore, $\formallimit{\frac{1}{n}} = 0$ because $d(\frac{1}{n}, 0) \leq \varepsilon$ for all $n \geq N.$
\end{proof}

\begin{definition}{Reciprocation of real numbers}{reciprocation}
    A sequence $(a_n)_{n=0}^\infty$ of rational numbers is said to be bounded away from zero if and only if
    there exists a rational number $c > 0$ such that $|a_n| \geq c$ for all $n \geq 1.$
    
    Let $x$ be a non-zero real number. 
    Let $(a_n)_{n=0}^\infty$ be a Cauchy sequence bounded away from zero such that $x = \formallimit{a_n}$ (such a sequence exists by Lemma 5.3.14).
    Then, we define the reciprocal $x^{-1}$ by the formula $x^{-1} := \formallimit{a_n^{-1}}$.
    (From Lemma 5.3.15, we know that $x^{-1}$ is a real number.)

    Finally, we define division $x.y$ of two real numbers $x, y$, provided $y$ is non-zero by the formula:
    $$x/y := x \times y^{-1}.$$
\end{definition}

\begin{problem}{Multiplicative Inverse}
    Show that $xx^{-1} = x^{-1}x= 1$.
\end{problem}

\begin{proof}
    Notice:
    \begin{align*}
        xx^{-1} &= \formallimit{a_n} \times \formallimit{a_n^{-1}} &&\text{by definition of $x$ and $x^{-1}$}\\
        &= \formallimit{a_n \times a_n^{-1}} &&\text{by definition of multiplication over the reals} \\
        &= \formallimit{1} &&\text{by the multiplicative inverse of the rationals}\\
        &= 1 &&\text{by the embedding of $\mathbb{Q}$ to $\mathbb{R}$}
    \end{align*}
    Because multiplication is commutative, we also have $x^{-1} x = 1.$
\end{proof}